\scanpage{287}
(1) $S$ anneau \inserted{$\mathbb Z$-}gradué \inserted{comm} \uline{à degrés $\geq0$}, $S=\sum_{n\geq0}S_n^{\unsure{b|h}}$

\deleted{$\unsure{S_0=k}$}

(2) $k=S_0^{\unsure{b|h}}$

(3) $\operatorname{Modgr.d}(S)=\mathcal M$

(4) $\underline t$ foncteur translation
\[
(4.1)\quad\left\{\begin{array}{l}
\underline t:\mathcal M\longrightarrow\mathcal M\\
\underline t(M)_i=M_{i+1}
\end{array}\right.
\]
\margin{c'est un \uline{automorphisme} de $\mathcal M$}
\[
(4.2)\qquad M[i]=\underline t^{\,i}(M).
\]
\deleted{\unsure{Il en va}}
Pour tout entier $N$, on considère les deux sous-catégories \inserted{str.} pleines \inserted{$\mathcal M_N,\mathcal M^N$} de $\mathcal M$ \unsure{définies} par
\[
(5)\quad\left\{\begin{array}{l}
\operatorname{Ob}\mathcal M_N=\{P\in\operatorname{Ob}\mathcal M\mid P_i=0\quad\forall i<N\}\\
\operatorname{Ob}\mathcal M^N=\{P\in\operatorname{Ob}\mathcal M\mid P_i=0\quad\forall i\unsure{>|\geq}N\}
\end{array}\right.
\]
\uline{NB} Les $\mathcal M_N$ ($\mathcal M^N$) décroissent (croissent) pour $N$ croissant. On a
\[
(6)\quad\left\{\begin{array}{l}
t^i(\mathcal M_N)=\mathcal M_{N-i}\\
t^i(\mathcal M^N]=\mathcal M^{N-i}_{\deleted{\unsure{N-i}}}
\end{array}\right.
\]
\deleted{\unsure{donc}}
En particulier \deleted{\unsure{et}}
\[
(6.1)\quad\left\{\begin{array}{l}
\deleted{\unsure{t}}\mathcal M_N=t^{-N}\mathcal M_0\quad,\\
\mathcal M^N=t^{-N}\mathcal M^0
\end{array}\right.
\]
\inserted{Dans} La catégorie $\mathcal M$, munie déjà de l'\unsure{autom.} $t$, la donnée des familles $(\mathcal M_N)$ et $(\mathcal M^N)$ avec leurs inclusions \unsure{définies}
\[
(7)\qquad\mathcal M_N\supset\mathcal M_{N'}\quad\text{et}\quad\mathcal M^N\subset\mathcal M^{N'}\quad\text{si}\quad N\leq N'
\]
équivaut à celle \deleted{\unsure{des}} du couple
\[
(\unsure{8|7})\qquad\mathcal M_0,\mathcal M^0\quad\text{avec}\quad t\mathcal M_0\unsure{\subset|\supset}\mathcal M_0,\quad
 t\mathcal M^0_{\deleted{\unsure{0}}}\supset\mathcal M^0_{\deleted{\unsure{0}}}
\]
\unsure{On va voir} d'ailleurs comment $\mathcal M_N,\mathcal M^{N-1}$ (\deleted{\unsure{ou encore $\mathcal M_N$, $\mathcal M^N$}}) se déterminent mutuellement -- donc $\mathcal M_0\simeq\mathcal M^0$ permet de \unsure{retrouver} tous \deleted{\unsure{les}} $\mathcal M_N,\mathcal M^N$.

\margin{288}
